\lecture{12}{Ускорение при сильной выпуклости}

\sourcevideo
    {https://www.youtube.com/playlist?list=PLOzRYVm0a65fhKDS8cYIC7YCuVGJ4FOJx}
    {Week 4: Lecture 12: Acceleration under Strong Convexity}

Рассмотрим градиентный спуск с шагом \(1/L\):
\[
    x^{k+1}=x^k-\frac1L\nabla f(x^k).
\]
Для выпуклой \(L\)-гладкой функции он имеет скорость
\(\mathcal O(1/k)\). Если функция дополнительно удовлетворяет
неравенству Поляка--Лоясиевича, та же итерация сходится линейно.
Таким образом, в этой лекции ускоряется не алгоритм, а теоретическая
оценка за счёт более сильной геометрии целевой функции.

\section{Лемма об убывании}

Рассматривается задача
\[
    \minimize_{x\in\RR^n} f(x),
\]
где минимум достигается:
\[
    x^\star\in\operatorname*{argmin}_{x\in\RR^n}f(x),
    \qquad
    f^\star=f(x^\star).
\]
Гладкость с константой \(L\) означает
\[
    \norm{\nabla f(x)-\nabla f(y)}
    \le
    L\norm{x-y}.
\]

\begin{lemma}[Квадратичная верхняя оценка]
Если \(f\) является \(L\)-гладкой, то для любых \(x,y\in\RR^n\)
\[
    f(y)
    \le
    f(x)
    +
    \inner{\nabla f(x)}{y-x}
    +
    \frac L2\norm{y-x}^2.
\]
\end{lemma}

Подставляя
\[
    y=x-\frac1L\nabla f(x),
\]
получаем
\[
\begin{aligned}
    f\left(x-\frac1L\nabla f(x)\right)
    &\le
    f(x)
    -\frac1L\norm{\nabla f(x)}^2
    +\frac1{2L}\norm{\nabla f(x)}^2
    \\
    &=
    f(x)-\frac1{2L}\norm{\nabla f(x)}^2.
\end{aligned}
\]
Поэтому для градиентного спуска
\[
    f(x^{k+1})
    \le
    f(x^k)
    -
    \frac1{2L}\norm{\nabla f(x^k)}^2,
\]
и последовательность значений функции невозрастает.

\section{Выпуклый режим}

Для получения явной скорости используется более сильная форма
предыдущей оценки.

\begin{lemma}[Трёхточечное неравенство]
Пусть \(f\) выпукла и \(L\)-гладка, а
\[
    x^{k+1}=x^k-\frac1L\nabla f(x^k).
\]
Тогда для любого \(x\in\RR^n\)
\[
\begin{aligned}
    f(x^{k+1})
    \le
    f(x)
    +
    \frac L2
    \left(
        \norm{x^k-x}^2
        -
        \norm{x^{k+1}-x}^2
    \right).
\end{aligned}
\]
\end{lemma}

\begin{proof}
Из гладкости
\[
\begin{aligned}
    f(x^{k+1})
    \le{}&
    f(x^k)
    +
    \inner{\nabla f(x^k)}{x^{k+1}-x^k}
    \\
    &+
    \frac L2\norm{x^{k+1}-x^k}^2.
\end{aligned}
\]
По выпуклости
\[
    f(x^k)
    \le
    f(x)
    +
    \inner{\nabla f(x^k)}{x^k-x}.
\]
Складывая неравенства и используя
\[
    \nabla f(x^k)=L(x^k-x^{k+1}),
\]
получаем
\[
\begin{aligned}
    f(x^{k+1})
    \le{}&
    f(x)
    +
    L\inner{x^k-x^{k+1}}{x^{k+1}-x}
    \\
    &+
    \frac L2\norm{x^{k+1}-x^k}^2.
\end{aligned}
\]
Теперь применяется тождество
\[
    2\inner{a-b}{b-c}
    =
    \norm{a-c}^2
    -
    \norm{a-b}^2
    -
    \norm{b-c}^2.
\]
\end{proof}

Полагая \(x=x^\star\), имеем
\[
\begin{aligned}
    f(x^{k+1})-f^\star
    \le
    \frac L2
    \left(
        \norm{x^k-x^\star}^2
        -
        \norm{x^{k+1}-x^\star}^2
    \right).
\end{aligned}
\]
Суммирование по \(k=0,\ldots,K\) даёт
\[
    \sum_{k=0}^{K}
    \bigl(f(x^{k+1})-f^\star\bigr)
    \le
    \frac L2\norm{x^0-x^\star}^2.
\]
Поскольку \(f(x^k)\) невозрастает,
\[
\begin{aligned}
    (K+1)
    \bigl(f(x^{K+1})-f^\star\bigr)
    \le
    \frac L2\norm{x^0-x^\star}^2.
\end{aligned}
\]

\begin{theorem}[Скорость для выпуклой гладкой функции]
Пусть \(f\) выпукла и \(L\)-гладка. Тогда
\[
    f(x^{K+1})-f^\star
    \le
    \frac{L\norm{x^0-x^\star}^2}{2(K+1)}.
\]
\end{theorem}

Следовательно,
\[
    f(x^k)-f^\star
    =
    \mathcal O\!\left(\frac1k\right),
    \qquad
    k
    =
    \mathcal O\!\left(\frac1\varepsilon\right)
\]
для достижения ошибки не более \(\varepsilon\). Термин
«сублинейная скорость» относится к верхней оценке \(C/k\); он не
означает, что фактическое отношение соседних ошибок обязано равняться
\(k/(k+1)\).

\section{PL-условие и линейная сходимость}

\begin{definition}[Неравенство Поляка--Лоясиевича]
Дифференцируемая функция \(f\) удовлетворяет PL-неравенству с
константой \(\mu>0\), если
\[
    \frac12\norm{\nabla f(x)}^2
    \ge
    \mu\bigl(f(x)-f^\star\bigr)
\]
для всех \(x\).
\end{definition}

Всякая дифференцируемая \(\mu\)-сильно выпуклая функция
удовлетворяет PL-неравенству с той же константой, но дальнейшее
доказательство использует только PL-свойство.

Из леммы об убывании
\[
    f(x^{k+1})
    \le
    f(x^k)
    -
    \frac1{2L}\norm{\nabla f(x^k)}^2.
\]
По PL-неравенству
\[
    \frac12\norm{\nabla f(x^k)}^2
    \ge
    \mu\bigl(f(x^k)-f^\star\bigr).
\]
Следовательно,
\[
\begin{aligned}
    f(x^{k+1})-f^\star
    \le
    \left(1-\frac\mu L\right)
    \bigl(f(x^k)-f^\star\bigr).
\end{aligned}
\]

\begin{theorem}[Линейная сходимость при PL-условии]
Пусть \(f\) является \(L\)-гладкой и удовлетворяет
PL-неравенству с константой \(\mu>0\). Тогда
\[
    f(x^k)-f^\star
    \le
    \left(1-\frac\mu L\right)^k
    \bigl(f(x^0)-f^\star\bigr).
\]
\end{theorem}

Для \(L\)-гладкой PL-функции можно считать \(0<\mu\le L\), поэтому
коэффициент \(1-\mu/L\) принадлежит \([0,1)\). Ошибка уменьшается
геометрически: если она ещё велика, PL-неравенство не позволяет норме
градиента быть слишком малой.

Так как \(\nabla f(x^\star)=0\), из гладкости следует
\[
    f(x^0)-f^\star
    \le
    \frac L2\norm{x^0-x^\star}^2.
\]
Поэтому
\[
    f(x^k)-f^\star
    \le
    \frac L2
    \left(1-\frac\mu L\right)^k
    \norm{x^0-x^\star}^2.
\]
При нескольких минимизаторах \(x^\star\) можно выбирать произвольно;
однозначность гарантируется только сильной выпуклостью.

\section{Число обусловленности и сложность}

Для \(L\)-гладкой и \(\mu\)-сильно выпуклой функции вводится число
обусловленности
\[
    \kappa=\frac L\mu\ge1.
\]
Тогда коэффициент сходимости равен
\[
    1-\frac\mu L=1-\frac1\kappa.
\]
При \(\kappa\approx1\) кривизна функции почти одинакова во всех
направлениях. При \(\kappa\gg1\) линии уровня вытянуты, а коэффициент
сжатия близок к единице.

Для квадратичной функции
\[
    f(x)=\frac12x^{\mathsf T}Qx+c^{\mathsf T}x,
    \qquad
    Q\succ0,
\]
можно взять
\[
    L=\lambda_{\max}(Q),
    \qquad
    \mu=\lambda_{\min}(Q),
\]
поэтому
\[
    \kappa
    =
    \frac{\lambda_{\max}(Q)}{\lambda_{\min}(Q)}.
\]

Обозначим
\[
    \Delta_k=f(x^k)-f^\star.
\]
Из неравенства \(1-t\le e^{-t}\) получаем
\[
    \Delta_k
    \le
    \exp\left(-\frac\mu Lk\right)\Delta_0.
\]
Чтобы обеспечить \(\Delta_k\le\varepsilon\), достаточно
\[
    k
    \ge
    \frac L\mu
    \log\left(\frac{\Delta_0}{\varepsilon}\right).
\]
Следовательно,
\[
    k
    =
    \mathcal O\!\left(
        \kappa\log\frac1\varepsilon
    \right)
\]
при фиксированной начальной ошибке. В выпуклом случае зависимость была
\(\mathcal O(1/\varepsilon)\); PL-свойство заменяет её логарифмической
зависимостью от требуемой точности.

\section{PL-класс шире сильной выпуклости}

Сильная выпуклость влечёт PL-неравенство, но обратная импликация
неверна даже для выпуклых гладких функций. Рассмотрим
\[
    f(x_1,x_2)=\frac12x_1^2.
\]
Тогда
\[
    \nabla f(x_1,x_2)
    =
    \begin{pmatrix}
        x_1\\
        0
    \end{pmatrix},
    \qquad
    f^\star=0,
\]
и
\[
    \frac12\norm{\nabla f(x)}^2
    =
    \frac12x_1^2
    =
    f(x)-f^\star.
\]
Следовательно, функция удовлетворяет PL-неравенству с \(\mu=1\).
Однако она постоянна вдоль направления \((0,1)^{\mathsf T}\), не
является сильно выпуклой и имеет множество минимумов
\[
    \operatorname*{argmin}f
    =
    \setgiven{(0,x_2)}{x_2\in\RR}.
\]
Таким образом,
\[
    \text{сильная выпуклость}
    \Longrightarrow
    \text{PL-неравенство},
\]
но не наоборот.

\section{Интерпретация и ограничения}

В обоих режимах используется одна и та же итерация
\[
    x^{k+1}=x^k-\frac1L\nabla f(x^k).
\]
Для общего выпуклого класса градиент может быть малым без
количественной связи с ошибкой по значению, поэтому получается оценка
порядка \(1/k\). PL-неравенство устанавливает связь
\[
    \norm{\nabla f(x)}^2
    \ge
    2\mu\bigl(f(x)-f^\star\bigr),
\]
из-за которой каждая итерация уменьшает текущую ошибку как минимум в
фиксированное число раз.

Шаг \(1/L\) требует знания или оценки константы гладкости. Для
предсказания линейной скорости нужна также оценка \(\mu\). Если
\(\mu\ll L\), то большое число обусловленности делает даже линейную
сходимость медленной. Одна итерация использует один полный градиент и
хранит текущий вектор параметров. При стохастическом или
распределённом градиенте появляются дополнительные ошибки, связанные
с дисперсией и коммуникациями, поэтому оценки этой лекции напрямую к
ним не переносятся.
